% !TeX program = xelatex
% =====================================================================
%  示例简历 04 · 数据开发工程师 · 社招（无证件照版式）
%  写法说明：社招写"负责范围 + 治理成效"，数字按规模、成本、事故数来写，
%  不写单接口延迟；工作经历比项目经历更靠前。
%  【虚构声明】人物用通用占位名（赵六），学校与公司一律以“示例”开头，均为虚构。
%  奖项、指标同样是编造的，不指向任何真实个人、学校或公司。
%  编译：xelatex 04-data-engineer.tex   （或在本仓库 make examples）
% =====================================================================
\documentclass[11pt,a4paper]{article}

\usepackage[UTF8,fontset=none]{ctex}
\usepackage{fontspec}
\usepackage[left=1.15cm,right=1.15cm,top=0.85cm,bottom=0.8cm]{geometry}
\usepackage{enumitem}
\usepackage{tabularx}
\usepackage{titlesec}
\usepackage{xcolor}
\usepackage{graphicx}
\usepackage{microtype}
\usepackage{xurl}
\usepackage{hyperref}

\defaultfontfeatures{
    Extension=.otf,
    UprightFont=*-regular,
    BoldFont=*-bold,
    ItalicFont=*-italic,
    BoldItalicFont=*-bolditalic,
    Numbers=Lining}
\setmainfont{texgyrepagella}
\setsansfont{texgyrepagella}
\defaultfontfeatures{}   % 清空默认项，下面的中文字体各自显式指定
% 正文中文用楷体；楷体无粗体，加粗与标题用黑体，混排更清晰
\IfFontExistsTF{KaiTi}{%
    \IfFontExistsTF{SimHei}{%
        \setCJKmainfont{KaiTi}[BoldFont=SimHei]%
    }{%
        \setCJKmainfont{KaiTi}[AutoFakeBold=2.5]%
    }%
}{%
    \IfFontExistsTF{FandolKai-Regular.otf}{%
        \setCJKmainfont{FandolKai-Regular.otf}[BoldFont=FandolHei-Regular.otf]%
        \setCJKsansfont{FandolHei-Regular.otf}%
    }{%
        \setCJKmainfont{Noto Serif CJK SC}[BoldFont=Noto Sans CJK SC]%
        \setCJKsansfont{Noto Sans CJK SC}%
    }%
}

\definecolor{theme}{HTML}{1A2B3C}       % 主色：沉稳炭青
\definecolor{textmain}{HTML}{222222}    % 正文
\definecolor{textmuted}{HTML}{5A5A5A}   % 次级信息
\definecolor{linecolor}{HTML}{D0D7DE}   % 分割线

\color{textmain}
\pagestyle{empty}
\setlength{\parindent}{0pt}
\setlength{\parskip}{0pt}
\linespread{1.15}
\emergencystretch=2em
\tolerance=3000

% 带下划线的链接（URL 短，不会跨行）
\newcommand{\ulink}[2]{\href{#1}{\color{theme}\underline{#2}}}
\newcommand{\skillitem}[2]{\item \textbf{#1}：\,#2}

% 经历抬头：标题靠左、身份/方向居中、时间靠右（等间距分布，标题过长只会报 Overfull，不会压字）
\newcommand{\entry}[3]{%
    \par\addvspace{2.9pt}%
    \noindent
    \hbox to \linewidth{%
        {\bfseries\fontsize{11.5}{14}\selectfont\color{theme}#1}%
        \hfil{\small\color{textmuted}#2}%
        \hfil{\small\color{textmuted}#3}%
    }\par\nobreak\vspace{1pt}%
}
% 经历下方的一行说明（实习业务方向 / 项目简介）
\newcommand{\subline}[1]{{\small\color{textmuted}#1}\par\nobreak\vspace{1.5pt}}

\titleformat{\section}
    {\large\bfseries\color{theme}}{}{0pt}{}
    [{\vspace{1.5pt}\color{linecolor}\hrule height 0.7pt}]
\titlespacing*{\section}{0pt}{8.4pt}{3pt}

\setlist[itemize]{
    leftmargin=1.15em,
    labelsep=0.45em,
    itemsep=2.9pt,
    parsep=0pt,
    topsep=1.8pt,
    label={\small\color{theme}\textbullet}
}

\hypersetup{
    colorlinks=true,
    urlcolor=theme,
    linkcolor=theme,
    pdfauthor={你的姓名},
    pdftitle={你的姓名 · 个人简历}
}

\begin{document}
\hypersetup{pdfauthor={赵六}, pdftitle={赵六 · 个人简历}}

{\bfseries\fontsize{25}{28}\selectfont\color{theme}赵六}%
\hspace{0.85em}{\fontsize{13.5}{16}\selectfont\color{theme}数据开发工程师}\par
\vspace{4pt}
{\small\color{textmuted}男 · 26 岁 · 本科 · 上海 · 4 年数据开发经验}\par
\vspace{3pt}
{\small
    136-0000-0004 \quad
    zhaoliu@example.com \quad
    \ulink{https://github.com/example}{github.com/example}
}

\section{教育背景}
\entry{示例大学}{软件工程 · 本科 · GPA 3.6 / 4.0}{2019.09 — 2023.06}
\subline{主修课程：数据库系统、数据结构与算法、操作系统、数据挖掘；校级二等奖学金两次。}

\section{工作经历}
\entry{示例信息}{高级数据开发工程师}{2025.04 — 至今}
\subline{实时数仓与指标平台，服务 6 条业务线的日报与实时看板。}
\begin{itemize}
    \item 用 \textbf{Flink + Kafka 双流 join} 替掉原来的小时级批处理，指标从"隔小时看"变成分钟级可见，日处理量在十亿条量级。
    \item 做了冷热分层、小文件合并和分区裁剪，一年下来存储成本降了三分之一多，核心报表查询也快了一倍多。
    \item 牵头整理指标字典与血缘，重复开发的指标少了一大半，口径争议的工单也不再每周都来。
    \item 把重复的实时链路抽成配置化模板，新业务接入从两周缩到三天，不用每次重写一遍消费和落库逻辑。
    \item 主导数据需求评审与方案评审，带两名新同学完成模块交付，需求返工的情况少了很多。
\end{itemize}

\entry{示例出行}{数据开发工程师}{2023.07 — 2025.03}
\subline{用户增长方向的数据链路与报表体系，对接产品与运营需求。}
\begin{itemize}
    \item 把两百多个散落脚本迁到 DAG 调度与依赖管理，任务失败率从 8\% 降到 1.2\%，夜里被叫起来重跑的次数少了很多。
    \item 上线规则校验与延迟告警，覆盖八十多张核心表，一个季度的数据事故从五起降到零到一起。
    \item 把手工 Excel 报表迁移成自助看板，运营同学每周省下十几个小时的人工统计时间。
\end{itemize}

\section{项目经历}
\entry{数据质量平台 — 规则引擎与告警}{项目负责人 · Python / Flink / ClickHouse}{2024.06 — 2025.02}
\subline{面向全公司数据表的可配置校验平台：规则 DSL、采样执行、分级告警与回滚。}
\begin{itemize}
    \item 设计了规则 DSL 与执行计划，支持空值、枚举、波动、主键唯一等十二类校验，业务方按模板填配置即可接入。
    \item 告警按影响面分级并做聚合抑制，值班同学从一天几十条告警变成只看真正要紧的几条。
    \item 规则配置与执行解耦后，六成左右的校验规则业务同学能自己配，不用每次排期等开发。
    \item 上线八个月覆盖 6 条业务线，提前发现二十多起数据异常，其中两次本来会直接发到对外报表上。
\end{itemize}

\entry{实时标签平台 — 用户画像标签}{核心开发 · Flink / HBase / Redis}{2023.10 — 2024.05}
\subline{为增长与推荐提供用户标签，离线与实时双链路共出一百多个标签。}
\begin{itemize}
    \item 离线回溯与实时增量共用一套口径，标签和离线结果对不上时能直接查出是哪条链路的问题。
    \item 查询侧用 Redis 冷热分层 + 批量接口，推荐每次请求都要查标签，高峰期也没出现过超时。
    \item 消费积压时能自动扩容，配合断点续传，流量翻几倍也没有丢过数据。
\end{itemize}

\section{专业技能}
\begin{itemize}
    \skillitem{编程语言}{Python（主力）、SQL、Scala、Java、Shell}
    \skillitem{大数据组件}{Hadoop / Hive、Spark、Flink、Kafka、ClickHouse、Doris}
    \skillitem{数仓与治理}{维度建模、离线 / 实时数仓、指标字典、数据质量与血缘}
    \skillitem{工程与平台}{Airflow / DolphinScheduler、Docker、K8s 基础、阿里云 / 腾讯云组件}
    \skillitem{协作与其他}{需求与方案评审、数据产品对接、技术文档与新人带教}
\end{itemize}

\end{document}
